\section{Threats to validity}
\label{sec:threats}

The design removes some sources of bias and leaves others in place.
This section lists the ones it cannot remove, in roughly the order of how much I expect them to matter.

\paragraph{Self-selection within solvers.}
Participants choose both their parents and their models.
Solver strata remove stable traits, such as skill or compute budget, from the primary contrast.
They do not remove choices that vary within a solver.
A solver who reaches for a different model when a problem looks harder, or who builds on another solver's work only after a promising idea appears, changes the comparison in ways the covariates may not capture.
The time spline and the current-best indicator absorb part of this, and the rest remains.

\paragraph{Self-declared labels.}
The \field{Model:} line is what the solver typed, and the platform does not verify it.
In the audit, 93.2\% of non-silent notes agreed with the declared model and 94.6\% with the declared family, below the 95\% bar set in advance (Section~\ref{sec:labels}).
Most disagreements were multi-vendor work under a single-vendor label, which would move some cross-family rows into the same-family arm.
Uniform label noise of this kind tends to pull an odds ratio toward 1, so it bears on how to read an equivalence verdict.
Sensitivity analyses 1, 4 and 6 change how labels and harnesses enter the model.
The audit itself was scored by models from one of the vendors being labelled.

\paragraph{Parent inference.}
Declared attribution exists for only some rows, and the tree-distance parent is an inference from code similarity.
A solver can build on an idea without copying its code, and two solvers can reach similar code independently.
AT4 measures how often the inferred cell matches the cell implied by declared or text-named parents (Section~\ref{sec:parents}).
Sensitivity analysis 2 replaces every declared parent with the tree parent.
Rows whose tree is far from every candidate are treated as novel and leave the primary sample, which changes who is compared with whom.

\paragraph{Missing attempts.}
The API holds only submitted work.
Local experiments, discarded branches and private agent sessions are absent, a point the ECDSA.Fail paper also makes about its own archive \citep[\S6]{long2026ecdsafail}.
Y1 is therefore promotion per evaluated submission, never per attempt.
If solvers submit cross-family work only when it looks promising, the per-submission rate overstates its per-attempt rate, and the same holds for same-family work.

\paragraph{Missing trees.}
On ECDSA.Fail, commit availability differs between accepted and rejected rows in the early weeks.
The confirmatory window starts only once tree recovery is balanced between promoted and other rows, and AT15 checks this.
Rows before the window are excluded, so results describe the later part of each benchmark.

\paragraph{The denominator.}
The line between an evaluated failure and an infrastructure failure was drawn from reason text.
Ambiguous reasons are kept in the primary analysis and dropped in sensitivity analysis 9.

\paragraph{Dependence between rows.}
Resubmissions of one idea, shared parents and repeated nonce draws make rows within a benchmark dependent in ways the solver strata do not fully capture.
Sensitivity analysis 3 drops rows whose note reports only a nonce or validation-draw change.
Sibling benchmarks share tasks or repositories, and sensitivity analysis 8 keeps one member of each pair.

\paragraph{One platform.}
Every benchmark runs on Yukon, with one submission format and one promotion mechanism.
Promotion thresholds differ between benchmarks, and on at least one benchmark the threshold appears to have changed part-way through (Section~\ref{sec:outcomes}).
Results may not carry over to settings where participants coordinate differently or share work outside a leaderboard.

\paragraph{The motivating comparison.}
The gate counts in the 16 October post come from two searches with different interfaces and accounting conventions (Section~\ref{sec:intro}).
The study does not use that comparison as evidence, and its result says nothing about which search found the better circuit.

\paragraph{What was seen before the lock.}
Crude analyses on ECDSA.Fail were computed before the lock (Section~\ref{sec:peek}).
The parent rule, sample and model were fixed in writing before outcomes were joined, and the lock binds the code to that text.
A reader can still discount ECDSA.Fail's own estimate on these grounds, and the pooled primary is less exposed because the other benchmarks were not split by family before the lock.

\paragraph{Conflict of interest.}
I am a participant on ECDSA.Fail and a listed author of the paper whose claims this study tests (Section~\ref{sec:coi}).
Sensitivity analysis 7 removes my rows, and the full pipeline is public.

\paragraph{Reviewer independence.}
The code and the label audit were reviewed by Claude models, which share weights with the agents that wrote the code.
Agreement among them is one opinion, and the design has not yet had an outside statistical review.
%TODO Name any outside reviewer here after M6, or keep this sentence.
